%% Tier-2 proof compressions for v1.1 Part III, source lines 3022--3926.
%% Every item corresponds to one complete source proof and targets 40--60%
%% of that proof's detex/[A-Za-z]+ count.

\begin{annexitem}{annex:p2.2.syn.1}{thm:p02-fol-unique-readability}{Unique readability for first-order expressions}
\begin{proof}
For terms, disjoint tokens separate variables from compounds; the leading
function symbol and its fixed arity determine the argument count, and balanced
parentheses separate the comma-delimited arguments.  Induction on length gives
unique construction.  For formulas, equality and relation symbols with fixed
arities determine atoms.  The remaining initial tokens distinguish
\(\bottom\), negation, binary formulas, and quantified formulas.  A
balanced-depth scan finds the unique principal binary connective; removing a
unary symbol fixes its matrix.  Induction on formula length concludes.
\end{proof}
\end{annexitem}

\begin{annexitem}{annex:p2.2.syn.2}{thm:p02-fol-substitution-complexity}{Substitution preserves formula complexity}
\begin{proof}
Induct on complexity.  Atomic term replacement changes no logical constructor;
connective cases follow componentwise.  At a quantifier, substitution stops,
commutes, or first renames.  Each alternative preserves the quantifier, and
renaming preserves matrix complexity.
\end{proof}
\end{annexitem}

\begin{annexitem}{annex:p2.2.syn.3}{thm:p02-fol-free-variables-substitution}{Free variables under substitution}
\begin{proof}
Term induction proves both term claims, argumentwise in the function case.
For formulas use strong induction on complexity.  Atomic cases use those
claims; \(\bottom\) and connectives follow from their recursive clauses.

For variable replacement by fresh \(y\), no renaming clause fires.  At
\(Qw\,\psi\), substitution stops when \(w=x\); otherwise the induction
hypothesis passes through removal of \(w\), since \(w\ne y\), giving the
stated equality.  For the inclusion at \(Qy\,\psi\), the stopping and safe
commuting cases are immediate.  In the renaming case, first apply the equality
just proved to \(\psi\) under \([z/y]\), then the inclusion hypothesis to substitution
by \(t\).  Strong induction applies because
\(c(\psi[z/y])=c(\psi)<c(Qy\,\psi)\).  Freshness of \(z\) from
\(\operatorname{Var}(t)\cup\{x\}\), followed by removing the new binder,
gives the required subset.
\end{proof}
\end{annexitem}

\begin{annexitem}{annex:p2.2.syn.4}{thm:p02-fol-inert-substitution}{Substitution for a nonfree variable}
\begin{proof}
Use strong induction on complexity.  Atomic terms omit \(x\), and
\(\bottom\) is fixed; congruence handles connectives.  For \(Qy\,\psi\),
substitution stops if \(y=x\).  If \(y\ne x\) and
\(y\notin\operatorname{Var}(t)\), apply induction to \(\psi\).  Otherwise
the result is \(Qz((\psi[z/y])[t/x])\).  Freshness and the free-variable
lemma give \(x\notin\FV(\psi[z/y])\); complexity preservation licenses the
strong hypothesis.  Hence its matrix is alpha-equivalent to \(\psi[z/y]\).
Quantifier congruence, the generating fresh renaming from \(y\) to \(z\),
and transitivity yield \(Qz((\psi[z/y])[t/x])=_{\alpha} Qy\,\psi\).
\end{proof}
\end{annexitem}

\begin{annexitem}{annex:p2.2.syn.5}{thm:p02-fol-term-coincidence}{Term coincidence}
\begin{proof}
Induct on the term.  Agreement handles variables; argumentwise induction and
then \(f^{\mathcal M}\) handle compounds.
\end{proof}
\end{annexitem}

\begin{annexitem}{annex:p2.2.syn.6}{thm:p02-fol-term-substitution}{Term-substitution lemma}
\begin{proof}
Induct on \(r\).  Both variable cases are immediate; function terms follow
argumentwise from induction and the value clause.
\end{proof}
\end{annexitem}

\begin{annexitem}{annex:p2.2.syn.7}{thm:p02-fol-coincidence}{Formula coincidence}
\begin{proof}
Induct structurally, using term coincidence for atoms and the hypotheses for
connectives.  At \(Qx\,\psi\), equal updates make assignments agree on
\(\FV(\psi)\); quantify the resulting equivalence over every, respectively
some, \(a\).
\end{proof}
\end{annexitem}

\begin{annexitem}{annex:p2.2.syn.8}{thm:p02-fol-substitution}{Substitution lemma}
\begin{proof}
Induct on formula complexity.  Term substitution handles atoms and the
hypotheses handle connectives.  For \(Qy\,\psi\), if \(y=x\), the quantified
update overrides the earlier one.  If \(y\ne x\) and
\(y\notin\operatorname{Var}(t)\), term coincidence makes the value of \(t\)
unchanged under \(s[y\mapsto a]\); apply induction for each relevant \(a\).
Otherwise rename to fresh \(z\).  Freshness excludes \(z\) from \(t\) and
\(\FV(\psi)\).  Complexity preservation licenses induction on
\(\psi[z/y]\); induction on \(\psi\) for \([z/y]\), followed by formula
coincidence, equates \(Qz(\psi[z/y])\) with \(Qy\,\psi\).  The safe case
completes both quantifiers.
\end{proof}
\end{annexitem}

\begin{annexitem}{annex:p2.2.syn.9}{thm:p02-fol-alpha-invariance}{Alpha-invariance}
\begin{proof}
Let \(R\) relate formulas having identical satisfaction under every structure
and assignment.  Satisfaction makes \(R\) an equivalence and a congruence for
connectives and quantifiers.  For a generating renaming with
\(y\notin\operatorname{Var}(\varphi)\), the substitution lemma evaluates
\(\varphi[y/x]\) under \(s[y\mapsto a]\) as \(\varphi\) under
\(s[y\mapsto a][x\mapsto a]\).  Formula coincidence removes the irrelevant
\(y\)-update.  The quantifier clauses therefore relate
\(Qx\,\varphi\) and \(Qy(\varphi[y/x])\).  By the leastness defining
alpha-equivalence, \(=_{\alpha}\) is contained in \(R\).
\end{proof}
\end{annexitem}

\begin{annexitem}{annex:p2.2.syn.10}{thm:p02-fol-sentence-independence}{Assignment independence for sentences}
\begin{proof}
Coincidence applies; variables absent.
\end{proof}
\end{annexitem}

\begin{annexitem}{annex:p2.2.syn.11}{thm:p02-definitional-conservativity}{Conservativity of definitional extensions}
\begin{proof}
Theorem~\ref{thm:p02-reduct-invariance} gives invariance under reduct; all
formulas used for
consequence here are sentences.  Given \(\mathcal M\entails T\), formula
coincidence and
\(\FV(\varphi)\subseteq\{x_1,\ldots,x_n\}\) make truth of \(\varphi\)
after the displayed updates depend only on the tuple.  Let
\(R^{\mathcal M}_{\varphi}\subseteq(\dom\mathcal M)^n\) contain exactly
those tuples.

For any expansion \(\mathcal M'\), distinct quantified variables ensure that
unwinding \(\delta_R\) ranges over every tuple.  Its atomic side says that the
tuple lies in \(R^{\mathcal M'}\); reduct invariance says that its defining
side holds exactly for \(R^{\mathcal M}_{\varphi}\).  Thus
\(\mathcal M'\entails\delta_R\) iff
\(R^{\mathcal M'}=R^{\mathcal M}_{\varphi}\), proving
existence and uniqueness; old sentences keep \(T\) true.

If \(T\entails\sigma\), the reduct of any extended model satisfies
\(T\), hence \(\sigma\), and reduct invariance returns \(\sigma\) upstairs.
Conversely, uniquely expand any \(\mathcal M\entails T\); the assumed
extended consequence gives \(\sigma\), which reduct invariance returns to
\(\mathcal M\).  Therefore the old-signature consequences coincide.

For derivability, weakening gives the old-to-extended direction.  Conversely,
extended derivability gives extended consequence by soundness, hence old
consequence by the result just proved; under Choice, first-order completeness
then gives old derivability.
\end{proof}
\end{annexitem}
